\documentclass[./examen.tex]{subfiles}
\begin{document}
\begin{center}
\makebox[\textwidth]{Nombre y Apellido:\enspace\hrulefill}
\end{center}
\begin{questions}

 \question Un tanque de gas metano contiene 10 litros de gas a una presión de 10 atmósferas y una temperatura de 15°C. Si se coloca en una heladera y se enfría a 0°C, hasta que la presión cae a 5 atmósferas, ¿Cuál será presión del gas?

\question Un recipiente cerrado contiene 5 gramos de gas hidrógeno a una temperatura de 1°C y una presión de 1 atmósfera. Si se calienta el recipiente hasta que la temperatura alcanza los 27°C, ¿Cuál será la nueva presión del gas?

 \question Un globo de helio tiene un volumen de 5 litros a una presión de 1 atmósfera y una temperatura de 20°C. Si se disminuye la temperatura 0°C a presión constante , ¿Cuál será el nuevo volumen en el globo?
\question Una cámara de presión contiene 2 moles de gas nitrógeno a una temperatura de 25°C y una presión de 2 atmósferas. Si se reduce el volumen de la cámara a la mitad y se aumenta la temperatura a 55°C, ¿cuál será la nueva presión del gas?

\end{questions}
\end{document}